\section{Related Work}
\label{sec:related}

% TODO: full related work. Key references:
%   - BrainFlow: brainflow.readthedocs.io (cite)
%   - Muse SDK / Muse LSL: muselsl, muse-js
%   - AASM scoring: berry2017
%   - Dormio / dream incubation: maes2020dormio
%   - TMR for declarative memory: rasch2013
%   - TMR for creative insight: schreiner2024sleep
%   - Sleep staging with consumer EEG: kothe2019
%   - On-device LLM for EEG: relevant works TBD

The Muse family of consumer EEG devices has been used in a wide range of
research applications \citep{brainflow2022}. The Muse S, the device used
in the present work, has 4 unipolar frontal channels at 256 Hz. The
Muse-LSL and muse-js open-source projects provide Linux and browser-based
streaming for the Muse family.

The American Academy of Sleep Medicine (AASM) scoring standard
\citep{berry2017} defines 5 sleep stages (Wake, N1, N2, N3, REM) using
central and occipital EEG, EOG, and chin EMG. The Muse S's 4-channel
frontal configuration is insufficient for full AASM scoring; the
present work outputs a 4-class model that collapses N2/N3 and labels REM
as uncertain.

Dream incubation has been studied at the MIT Media Lab, with the Dormio
system \citep{maes2020dormio} demonstrating that hypnagogic-state audio
cues can incubate dream content. The present work extends this paradigm
to a fully automated, on-device pipeline with local LLM-based primer
generation and dream-report analysis.

Targeted Memory Reactivation (TMR) for declarative memory consolidation
is well-established \citep{rasch2013}. TMR for creative insight is
plausible but unproven; recent reviews \citep{schreiner2024sleep} note
the limited evidence base.

The 4-class sleep staging problem from consumer-grade EEG has been
addressed by several groups \citep{kothe2019}. The present work builds on
this literature with a Muse-S-specific adaptation that includes a
per-user alpha baseline calibration.
